\section{A bound on rare large values}
For a nonnegative random variable $X$ and a real threshold $a>0$, we have
\[
\Prob(X\ge a)\le\frac{\E X}{a}.
\]
To prove it, split the sum defining the expectation into the outcomes where $X\ge a$ and the rest:
\[
\E X=\sum_{\omega:\,X(\omega)\ge a}p(\omega)X(\omega)+\sum_{\omega:\,X(\omega)<a}p(\omega)X(\omega).
\]
In the first sum every value $X(\omega)$ is at least $a$, so that sum is at least $\sum_{X(\omega)\ge a}p(\omega)\,a=a\Prob(X\ge a)$. In the second sum every term is nonnegative, because $X$ is nonnegative. Hence $\E X\ge a\Prob(X\ge a)$, and dividing by the positive number $a$ gives the claim. This is \emph{Markov's inequality} in a finite setting. For example, if the average score in a class is $20$ and no score is negative, then at most a quarter of the students can have scored $80$ or more.

The nonnegativity assumption is not decorative. Large negative values could offset large positive ones in the expectation. For example, a variable equal to $100$ or $-100$ with equal probability has expectation zero but is at least $100$ with probability $1/2$.

\begin{principle}{Habit 8: choose a quantity whose average is easier than its best value}
When direct construction is difficult, define a finite experiment and count the desired or undesired features by indicators. Ask which averages can be computed exactly. Then state the logical conclusion carefully: one suitable outcome, a probability bound, or an actual efficient construction are different results.
\end{principle}
